%قبل از شروع مسئله ترجیحاً چیزی ننویسید. اما اگر خواستید بنویسید، در قالب لیستی (فقرات و شمارش) بنویسید.
%برای اضافه کردن نمره، بعدی از شروع‌مسئله، [x] را اضافه کنید و به جای x نمره مورد نظر را قرار دهید. برای زیرمسئله‌ها هم به همین صورت است.
%\شروع‌مسئله[10]
\شروع‌مسئله [14]
\\
شنگدباو و شنگدباوت 
\lr{"shengdebao \& shengdebaoette"}
 $n$
 تا رشته 
 $s_1, s_2,..., s_n$
 از حروف انگلیسی (که تعداد ثابت و 
 $O(1)$
 است) نوشته‌اند یک رشته دلخواه را «رمانتیک» می‌نامند اگر پیشوند لااقل یکی از رشته‌های 
 $s_1, ..., s_n$
 باشد.
 
 آن‌ها با هم یک بازی‌ای می‌کنند به اینصورت که رشته‌ای دارند و هر سری کاراکتر به انتهای آن اضافه می‌کنند. اگر در مرحله‌ای یکی از آن‌ها حرکتی کند و رشته حاصل رمانتیک نباشد او می‌بازد.  الگوریتمی از 
 $O(\sum_{i=1}^{n} |s_i|)$
 ارائه دهید که پیدا کند چه کسی استراتژی برد دارد.
%\متن‌سیاه{سوال عملی اول:}\\
%\متن‌سیاه{دراپ‌تست}\\

%\متن‌سیاه{ورودی:}\\
%در خط اول ورودی، ابتدا k و سپس m داده می‌شود.\\
%\متن‌سیاه{خروجی:}\\
%در تنها سطر خروجی، تعداد آیفون‌های مورد نیاز می‌آید. اگر در هیچ صورتی چنین کاری ممکن نبود ۱- چاپ می‌شود.\\
%\متن‌سیاه{محدودیت‌ها:}\\
%$1\leq k\leq 10^4$

%$m\leq10^9$
%\\
%\textbf{\متن‌سیاه{ورودی و خروجی نمونه}\\
%% تست اول
%\begin{latin}
% \begin{tabularx}{\textwidth}{|X | X|}
% \hline
% \textbf{stdin} &
% \textbf{stdout}\\
% \hline
% 100 100&
% 1\\
% \hline
% 
% \end{tabularx}
%\end{latin}
%
%% تست دوم
%
%\begin{latin}
% \begin{tabularx}{\textwidth}{|X | X|}
% \hline
% \textbf{stdin} &
% \textbf{stdout}\\
% \hline
% 340 20&
% 3\\
% \hline
% 
% \end{tabularx}
%\end{latin}}
% صورت اصلی مسئله را در این‌جا بنویسید.



\solution{
	همه رشته‌ها را در یک درخت trie می‌ریزیم اندازه این درخت ترای حداکثر برابر جمع طول رشته‌ها خواهد بود. 
	حالا مسئله را به شکل دیگری نگاه می‌کنیم رشته‌ای که می‌نویسیم معادل مسیری از ریشه به پایین است و وقتی می‌بازیم که به رأسی برسیم که دیگر بچه‌ای ندارد چرا که در اینصورت دیگر پیشوند هیچ رشته‌ دیگری نیست. بنابراین مسئله به این صورت است که درختی داریم و از ریشه دو نفر بازی می‌کنند یکی در میان هر کس رأس بعدی مسیر را از بین بچه‌های آخرین رأس انتخاب می‌کند کسی که به برگی برسد می‌بازد.
	
	به ازای هر رأس از درخت مثل $v$ مقدار 
	$dp_v$
	را برابر $Win$ می‌گذاریم اگر با شروع از آن رأس و حرکت به سمت پایین نفر اول ببرد و در غیر اینصورت مقدار آن‌را برابر $Lose$ می‌گذاریم. اگر $v$ برابر برگ بود مقدار آن حتماً $Lose$ خواهد بود. در غیر اینصورت در بین بچه‌های این رأس اگر رأسی مثل $u$ باشد که داشته‌باشیم $dp_u = Lose$ یعنی با وارد شدن به آن رأس در زیر مسئله‌ای نفر دوم، نفر اول است و می‌بازد؛ یعنی ما می‌بریم پس در این صورت به آن رأس رفته و $dp_v=  Win$ خواهد شد. اگر همه بچه‌های این رأس برابر $Win$ باشد یعنی هر کاری کنیم نفر دوم در زیرمسئله‌اش می‌برد لذا $dp_v = Lose$ خواهد بود. این کار را با استفاده از یک $dfs$ و بروزرسانی پس‌ترتیب درخت می‌توانیم انجام دهیم که در کل به‌اندازه حداکثر جمع طول رشته‌ها طول می‌کشد. در نهاین مقدار
	$dp_{root}$
	جواب نهایی مسئله را مشخص می‌کند.
	
}

%%%%%%%%%%%%%%%%%%%%%%%%%%%%%%%%%%%%%%%%%%%%%%%%%%%%%
% برخی دستورات مفید برای تمرین‌ها. %
%%%%%%%%%%%%%%%%%%%%%%%%%%%%%%%%%%%%%%%%%%%%%%%%%%%%%

% فرمول‌های نمونه:
%$f(x) = \Theta(x)$
%$T(x, y) \leq \log(x) + 2^y + \binom{2n}{n}$
% {\tt code\_function}


% لیست موارد بدون شماره‌گذاری:
%\شروع{فقرات}
% \فقره مورد اول 
% \شروع{فقرات}
% \فقره مورد اول داخلی 
% \شروع{فقرات}
% \فقره اولین مورد خیلی داخلی 
% \فقره دومین مورد خیلی داخلی 
% \پایان{فقرات}
% \فقره مورد دوم داخلی 
% \پایان{فقرات}
% \فقره مورد دوم 
%\پایان{فقرات}

% لیست موارد شماره‌گذاری شده:
%\شروع{شمارش}
% \فقره یک مورد 
% \فقره یک مورد دیگر 
% \فقره باز هم یک مورد دیگر
%\پایان{شمارش}

% برای نوشتن یک فرمول چند خطی که یک شماره بخوره. برای حالت‌های دیگر به سایت زیر مراجعه کنید:
% https://www.overleaf.com/learn/latex/Aligning_equations_with_amsmath
%\begin{align} 
% \begin{split}
% (x+y)^3 &= (x+y)^2(x+y) \\
% &= (x^2+2xy+y^2)(x+y) \\
% &= (x^3+2x^2y+xy^2) + (x^2y+2xy^2+y^3) \\
% &= x^3+3x^2y+3xy^2+y^3
% \end{split} 
%\end{align}

% نوشتن ماتریس و آرایه دوبعدی
%\begin{align}
% A = 
% \begin{bmatrix}
% A_{11} & A_{21} \\
% A_{21} & A_{22}
% \end{bmatrix}
%\end{align}

% اضافه کردن تصویر با شرح. برای اضافه کردن تصویر بدون شرح، فقط دستور درج‌تصویر را استفاده کنید.
%\شروع{شکل}[h!]
%\تنظیم‌ازوسط
% \درج‌تصویر[پهنا=0.5\پهنای‌متن]{img.jpg}
%\شرح{شرح تصویر}
%\برچسب{تصویر:۱}
%\پایان{شکل}
%\پایان{زیرمسائل}

%کد زیر را ببینید.
%%\lstinputlisting{source_filename.py}
%\begin{latin}
%\begin{lstlisting}
%# Import the modules
%import sys
%import random
%
%ans = True
%
%while ans:
%question = raw_input("Ask the magic 8 ball a question: (press enter to quit) ")
%
%answers = random.randint(1,8)
%
%if question == "":
%sys.exit()
%
%elif answers == 1:
%print "It is certain"
%
%elif answers == 2:
%print "Outlook good"
%
%elif answers == 3:
%print "You may rely on it"
%
%elif answers == 4:
%print "Ask again later"
%
%elif answers == 5:
%print "Concentrate and ask again"
%
%elif answers == 6:
%print "Reply hazy, try again"
%
%elif answers == 7:
%print "My reply is no"
%
%elif answers == 8:
%print "My sources say no"
%\end{lstlisting}
%\end{latin}